\chapter*{Summary}
\addcontentsline{toc}{chapter}{Summary}
\markboth{Summary}{}

This dossier turns the supervised project topic \emph{YOLO-Based Cross Feature
Fusion Mamba Net for Object Detection} into a research programme that can
support a journal article. It has three parts: an evidence base, a proposed
method, and an implementation plan.

\paragraph{Evidence base}
The four source papers were read in full, including their tables and
reference lists: MambaVision~\citep{hatamizadeh2025mambavision},
HPS-DETR~\citep{wang2025hpsdetr}, UAVDet~\citep{yang2026uavdet} and
CFMW~\citep{li2025cfmw}. They were then placed against work published up to
October 2026. Three findings shape everything that follows.
\begin{itemize}
  \item On still images, most of the accuracy gain these papers report comes
        from familiar devices. These are a stride-4 detection head, more
        capacity in early layers, and better multi-scale fusion. The state-space
        (Mamba) component itself adds between $+0.6$ and about $+3$ points of
        mAP, often at a large cost in speed or computation.
  \item Combining a YOLO detector with a Mamba fusion block for visible and
        thermal images is already well covered. Fusion-Mamba~\citep{dong2025fusionmamba},
        COMO~\citep{como2024}, WaveMamba~\citep{wavemamba2025},
        MS2Fusion~\citep{ms2fusion2025} and more than twenty
        Mamba-in-YOLO variants occupy that space. A further static fusion
        block would be judged incremental.
  \item No published method was found that carries a state-space memory
        causally across an unbounded stream of visible--thermal frames, at
        constant memory, for small objects. The nearest work,
        MambaST~\citep{mambast2024}, uses three-frame windows on one pedestrian
        dataset and does not evaluate streaming.
\end{itemize}

\paragraph{Proposed method}
\method{} (Cross-Feature Fusion Mamba Network) keeps the supervised topic
intact. It is a dual-stream YOLO detector with a stride-4 small-object head
and state-space fusion of visible and thermal features. Its streaming
variant, \methodS{}, adds a \emph{Temporal State Memory} that carries the
state-space hidden state from frame to frame, after correcting for camera
motion. Both variants rest on one observation. In a selective state-space
model the discretisation step $\Delta$ decides whether the state keeps its
memory or takes in new input. \method{} uses this single control for three
purposes: to discount an unreliable modality at a location, to carry objects
through gaps in a sensor feed, and to adjust for irregular frame intervals.
Five hypotheses with explicit falsification criteria are stated in
Section~\ref{sec:hypotheses}.

\paragraph{Implementation plan}
The system is built in four phases, each a usable product in its own right.
\begin{enumerate}
  \item Phase 1 detects objects in still images from curated datasets.
  \item Phase 2 detects and tracks objects in uploaded video.
  \item Phase 3 runs on live webcam and network camera streams.
  \item Phase 4, an optional extension, adds anonymous multi-camera
        analytics.
\end{enumerate}
The core benchmark is RGBT-Tiny~\citep{ying2025rgbttiny}. It contains 115
paired visible--thermal sequences, and more than 81\% of its objects are
smaller than $16\times16$ pixels. A notebook-driven workflow of sixteen
notebooks runs over a tested Python package. Training runs on Kaggle T4 GPUs,
while a local RTX~4050 laptop GPU is used for development and live
demonstrations.

\paragraph{Ethics and law}
The fictional Machine of \emph{Person of Interest} motivated the project, but
the system deliberately stops short of it. It does not identify people, keep
watch-lists or profile behaviour. Those functions are prohibited or
high-risk under the EU AI Act~\citep{euaiact2024} and constrained by India's
Digital Personal Data Protection Act~\citep{dpdp2023}. The system is a
perception layer that detects and tracks objects and anonymous people, with
faces blurred when footage is ingested.
